\PassOptionsToPackage{dvipsnames}{xcolor}

% draft会跳过文档中的所有图片。正式导出时需要删掉draft参数。
\documentclass[12pt, a4paper, oneside]{ctexart}

\usepackage{amsmath}
\usepackage{amssymb}
\usepackage{bm}
\usepackage{graphicx}
\usepackage{mathrsfs}
\usepackage{geometry}
\usepackage{framed}
\usepackage{caption}
\usepackage{listings}
\usepackage{fancyhdr}
\usepackage{booktabs}
\usepackage{makecell}
\usepackage{indentfirst}
\usepackage{authblk}
\usepackage{multicol}
% \usepackage{draftwatermark}       % 需要应用水印时取消注释
\usepackage{enumitem}
\usepackage[hidelinks]{hyperref}
\usepackage{tikz}
\usepackage{ulem}
\usetikzlibrary{positioning, shapes.geometric}

% 分栏线宽
\columnseprule=0.4pt

% 定制第二级无序列表的点样式
\setlist[itemize,2]{label=$\diamond$}

% 页边距
\geometry{a4paper, scale=0.8}

\pagestyle{fancy}

% 调整页眉高度，用于去除警告
\setlength{\headheight}{25pt}

\fancyhf{}      % 清空页眉页脚设置
\fancyhead[L] {
    % 工大计算机系logo
    \includegraphics[height=7mm]{./images/logo1.jpg}
}
\fancyhead[C]{《线性代数》复习}
\fancyhead[R]{\leftmark}    % 右侧页眉：当前章标题

% 页脚居中放置页码
\fancyfoot[C]{\thepage}

% 设置章节标题自动编号的格式
\ctexset{
  section/number=\chinese{section},
%   subsection/name={,},
%   subsection/number=\chinese{subsection}
}

% 行距。ctexart默认值为1.3
\linespread{1.5}

\lstset{
  language=pascal,
  basicstyle=\ttfamily,
  frame=single,
  numbers=left
}

% \SetWatermarkText{Eslzzyl整理}            % 设置水印内容
% \SetWatermarkLightness{0.9}             % 设置水印透明度 0-1
% \SetWatermarkScale{0.8}                   % 设置水印大小 0-1

\renewcommand{\headrulewidth}{1pt}  %页眉线宽，设为0可以去页眉线
\renewcommand{\footrulewidth}{1pt}  %脚注线的宽度

\definecolor{shadecolor}{RGB}{241, 241, 255}

\title{
    \includegraphics[width=0.3\textwidth]{images/hfut-badge.pdf}
    
    \vspace{20pt}
    《线性代数》总复习
}
\author{Eslzzyl}
\date{\today}

\newcounter{problemname}
\newenvironment{problem}{\begin{shaded}\stepcounter{problemname}\par\noindent\textbf{例题\arabic{problemname}. }}{\end{shaded}\par}
\newenvironment{solution}{\begin{shaded}\par\noindent\textbf{解答：}}{\end{shaded}\par}
% \newenvironment{solution}{\par\noindent\textbf{答案. }}{\par}
% \newenvironment{note}{\par\noindent\textbf{例题\arabic{problemname}的注记. }}{\\\par}
\newenvironment{note}{\par\noindent\textbf{注记. }}{\par}

\begin{document}

\maketitle
\newpage
\tableofcontents
\vspace{20pt}
% 如果在目录处有备注，可以写在这里。

\newpage

\section{行列式}

\subsection{行列式的性质}

行列式中的行与列具有同等的地位，行列式的性质凡是对行成立的对列也同样成立，反之亦然。

\begin{enumerate}
  \item 行列式和它的转置行列式相等。
  \item 对换行列式的两行（列），行列式变号。
  \begin{itemize}
    \item 推论：如果行列式有两行（列）完全相同，则此行列式等于零。（把这两行对换，得$D=-D$）
  \end{itemize}
  \item 行列式的某一行（列）中所有的元素都乘以同一个数$k$，等于用$k$乘以此行列式。
  \begin{itemize}
    \item 推论：行列式中某一行（列）的所有元素的公因子可以提到行列式的外面。
  \end{itemize}
  \item 行列式中如果有两行（列）元素成比例，则此行列式等于零。（由上面的两个推论得出）
  \item 若行列式的某一行（列）的元素都是两数之和，则行列式可以拆开。
  \item 把行列式的某行（列）乘以同一数然后加到另一行（列）对应的元素上去，行列式不变。
\end{enumerate}

\subsection{行列式按行（列）展开}

\begin{itemize}
  \item $M_{ij}$：余子式
  \item $A_{ij}$：代数余子式
\end{itemize}

\begin{enumerate}
  \item 行列式等于它的任一行（列）的各元素与其对应的代数余子式乘积之和。
  \item 行列式某一行（列）的元素与另一行（列）的对应元素的代数余子式乘积之和等于零。
\end{enumerate}

\subsection{一些特殊的行列式}

\subsubsection{范德蒙德（Vandermonde）行列式}

\begin{equation}
  D_n=\begin{vmatrix}
    1 & 1 & \cdots & 1 \\
    x_1 & x_2 & \cdots & x_n \\
    x_1^2 & x_2^2 & \cdots & x_n^2 \\
    \vdots & \vdots & & \vdots \\
    x_1^{n-1} & x_2^{n-1} & \cdots & x_n^{n-1}
  \end{vmatrix}=\prod_{n\geq i\geq j\geq 1} (x_i-y_i)
\end{equation}

\section{矩阵及其运算}

\subsection{线性方程组和矩阵}

\subsubsection{线性方程组}

\textbf{齐次线性方程组}：即常数项全为0的线性方程组：

\begin{equation}
  \left\{\begin{matrix}
    a_{11}x_1+a_{12}x_2+\cdots +a_{1n}x_n=0 \\
    a_{21}x_1+a_{22}x_2+\cdots +a_{2n}x_n=0 \\
    \cdots \\
    a_{m1}x_1+a_{m2}x_2+\cdots +a_{mn}x_n=0
   \end{matrix}\right.
\end{equation}

齐次线性方程组\textbf{一定}有零解（$x_1=x_2=\cdots =x_n=0$），\textbf{可能}有非零解。

\textbf{非齐次线性方程组}：即常数项不全为0的线性方程组：

\begin{equation}
  \left\{\begin{matrix}
    a_{11}x_1+a_{12}x_2+\cdots +a_{1n}x_n=b_1 \\
    a_{21}x_1+a_{22}x_2+\cdots +a_{2n}x_n=b_2 \\
    \cdots \\
    a_{m1}x_1+a_{m2}x_2+\cdots +a_{mn}x_n=b_2
   \end{matrix}\right.
\end{equation}

\subsubsection{矩阵}

行数与列数都为$n$的矩阵称为$n$阶矩阵或$n$阶方阵。$n$阶矩阵$\mathbf{A}$也记作$\mathbf{A}_n$。

也就是说，只要见到“$n$阶矩阵”或者$\mathbf{A}_n$这样的说法时，就表示矩阵是方阵。

\begin{itemize}
  \item 行矩阵（行向量）：

  \begin{equation*}
    \mathbf{A}=(a_1,a_2,\cdots,a_n)
  \end{equation*}

  \item {列矩阵（列向量）}
  
  \begin{equation*}
    \mathbf{B}=\begin{pmatrix}
      b_1 \\ 
      b_2 \\
      \cdots \\
      b_m 
     \end{pmatrix}
  \end{equation*}

\end{itemize}

行数列数都相等的矩阵称为\textbf{同型矩阵}。

零矩阵：$\mathbf{O}$

线性变换和矩阵之间存在一一对应关系。

除了主对角线之外的其他所有元素均为0的矩阵叫做\textbf{对角矩阵}，记作\footnote{同济线代中对角矩阵的符号是字母A去掉一横}

\begin{equation*}
  \mathbf{A}=\text{diag}(\lambda_1, \lambda_2,\cdots,\lambda_n)
\end{equation*}

单位矩阵$\mathbf{E}$是特殊的对角矩阵。

\subsection{矩阵的运算}

\subsubsection{矩阵乘法}

若有两个矩阵$\mathbf{A}$、$\mathbf{B}$满足$\mathbf{AB}=\mathbf{O}$，不能就此得出$\mathbf{A}=\mathbf{O}$或$\mathbf{B}=\mathbf{O}$的结论。

\subsubsection{矩阵转置}

矩阵的转置：$(\mathbf{AB})^T=\mathbf{B}^T \mathbf{A}^T$

\subsubsection{对称矩阵}

若方阵$\mathbf{A}$满足$\mathbf{A}^T=\mathbf{A}$，则称$\mathbf{A}$为对称矩阵。显然所有的对角矩阵都是对称矩阵。

\subsubsection{矩阵的行列式}

方阵$\mathbf{A}$的行列式：$\text{det}\ \mathbf{A}$或$|\mathbf{A}|$。

有性质$|\mathbf{AB}|=|\mathbf{A}|\cdot |\mathbf{B}|$。因此，一般来说$\mathbf{AB}\neq\mathbf{BA}$，但总有$|\mathbf{AB}|=|\mathbf{BA}|$。

\subsubsection{伴随矩阵}

由行列式$|\mathbf{A}|$的各个元素的代数余子式$A_{ij}$所构成的矩阵
\begin{equation*}
  \mathbf{A}^{*}=\begin{pmatrix}
    A_{11} & A_{21} & \cdots & A_{n1} \\
    A_{12} & A_{22} & \cdots & A_{n2} \\
    \vdots & \vdots & & \vdots \\
    A_{1n} & A_{2n} & \cdots & A_{nn}
  \end{pmatrix}
\end{equation*}

称为$\mathbf{A}$的伴随矩阵。注意这里有一个隐含的转置运算。

$2\times 2$矩阵的伴随矩阵：主对角线元素交换位置，副对角线元素变号。

有如下规律：
\begin{equation*}
  \mathbf{AA}^{*}=\mathbf{A}^{*}\mathbf{A}=|\mathbf{A}|\mathbf{E}
\end{equation*}

\subsection{逆矩阵}

\subsubsection{定义、性质和求法}

定义：对于$n$阶矩阵$\mathbf{A}$，如果有一个$n$阶矩阵$B$，使
\begin{equation*}
  \mathbf{AB}=\mathbf{BA}=\mathbf{E}
\end{equation*}

则说矩阵$\mathbf{A}$是可逆的，称$\mathbf{B}$是$\mathbf{A}$的逆矩阵。记作$\mathbf{B}=\mathbf{A}^{-1}$。此时$\mathbf{B}$是唯一的。

单元素矩阵的逆矩阵仍然是一个单元素矩阵，这唯一的元素等于原矩阵中唯一元素的倒数。

性质：
\begin{enumerate}
  \item 若$\mathbf{A}$可逆，则$|\mathbf{A}|\neq 0$，此时$\mathbf{A}$称为非奇异矩阵（反之则为奇异矩阵）。
  \item 若$|\mathbf{A}|\neq 0$，则$\mathbf{A}$可逆，且
  \begin{equation*}
    \mathbf{A}^{-1}=\frac{1}{|\mathbf{A}|}\mathbf{A}^{*}。
  \end{equation*}
\end{enumerate}

运算规律：
\begin{enumerate}
  \item $(\mathbf{A}^{-1})^{-1}=\mathbf{A}$
  \item $(\lambda \mathbf{A})^{-1}=\frac{1}{\lambda}\mathbf{A}^{-1}$
  \item $\mathbf{A}$、$\mathbf{B}$为同阶矩阵且均可逆，则有$(\mathbf{AB})^{-1}=\mathbf{B}^{-1}\mathbf{A}^{-1}$。
  \item 若$\mathbf{A}$可逆，则$\mathbf{A}^T$亦可逆，且$(\mathbf{A}^T)^{-1}=(\mathbf{A}^{-1})^T$。
\end{enumerate}

\subsubsection{逆矩阵的初步应用}

设$\varphi (x)=a_0+a_1 x+\cdots+a_m x^m$为$x$的$m$次多项式，$\mathbf{A}$为$n$阶矩阵，记
\begin{equation*}
  \varphi(\mathbf{A})=a_0 \mathbf{E}+a_1 \mathbf{A}+\cdots +a_m \mathbf{A}^m
\end{equation*}

$\varphi(\mathbf{A})$称为矩阵$\mathbf{A}$的$m$次多项式。

\subsection{克拉默法则}

如果线性方程组
\begin{equation*}
  \left\{\begin{matrix}
    a_{11}x_1+a_{12}x_2+\cdots +a_{1n}x_n=b_1 \\
    a_{21}x_1+a_{22}x_2+\cdots +a_{2n}x_n=b_2 \\
    \cdots \\
    a_{m1}x_1+a_{m2}x_2+\cdots +a_{mn}x_n=b_2
   \end{matrix}\right.
\end{equation*}

的系数矩阵的行列式不等于零，即
\begin{equation*}
  |\mathbf{A}|=\begin{vmatrix}
    a_{11} & \cdots & a_{1n} \\
    \vdots & & \vdots \\
    a_{n1} & \cdots & a_{nn}
  \end{vmatrix}\neq 0
\end{equation*}

那么，这方程组有唯一解
\begin{equation*}
  x_1=\frac{|\mathbf{A}_1|}{|\mathbf{A}|},\quad x_2=\frac{|\mathbf{A}_2|}{|\mathbf{A}|},\cdots, \quad x_n=\frac{|\mathbf{A}_n|}{|\mathbf{A}|}
\end{equation*}

其中$\mathbf{A}_j$是把系数矩阵$\mathbf{A}$中第$j$列的元素用方程组右端的常数项代替后所得到的$n$阶矩阵。

\subsection{矩阵分块法}

\subsubsection{分块矩阵的相加和数乘}

分块矩阵的相加、数乘是显然的。其中相加要求两个矩阵的行数、列数都相同，且采用相同的分块方法。

\subsubsection{分块矩阵相乘}

分块矩阵相乘：设$\mathbf{A}$为$m\times l$矩阵，$\mathbf{B}$为$l\times n$矩阵，分块成
\begin{equation*}
  \mathbf{A}=\begin{pmatrix}
    \mathbf{A}_{11} & \cdots & \mathbf{A}_{1t} \\
    \vdots & & \vdots \\
    \mathbf{A}_{s1} & \cdots & \mathbf{A}_{st}
  \end{pmatrix},\quad
  \mathbf{B}=\begin{pmatrix}
    \mathbf{B}_{11} & \cdots & \mathbf{B}_{1r} \\
    \vdots & & \vdots \\
    \mathbf{B}_{t1} & \cdots & \mathbf{B}_{tr}
  \end{pmatrix}
\end{equation*}

其中$\mathbf{A}_{i1},\mathbf{A}_{i2},\cdots,\mathbf{A}_{it}$的列数分别等于$\mathbf{B}_{1j},\mathbf{B}_{2j},\cdots,\mathbf{B}_{tj}$的行数，那么
\begin{equation*}
  \mathbf{AB}=\begin{pmatrix}
    \mathbf{C}_{11} & \cdots & \mathbf{C}_{1r} \\
    \vdots & & \vdots \\
    \mathbf{C}_{s1} & \cdots & \mathbf{B}_{sr}
  \end{pmatrix}
\end{equation*}

其中
\begin{equation*}
  \mathbf{C}_{ij}=\sum_{k=1}^t \mathbf{A}_{ik} \mathbf{B}_{kj}\quad (i=1,\cdots,s;\ j=1,\cdots,r).
\end{equation*}

\subsubsection{分块矩阵的转置}

分块矩阵的转置：矩阵本身要转置，各个子块自身也要转置。设
\begin{equation*}
  \mathbf{A}=\begin{pmatrix}
    \mathbf{A}_{11} & \cdots & \mathbf{A}_{1r} \\
    \vdots & & \vdots \\
    \mathbf{A}_{s1} & \cdots & \mathbf{A}_{sr}
  \end{pmatrix}
\end{equation*}

则
\begin{equation*}
  \mathbf{A}^T=\begin{pmatrix}
    \mathbf{A}_{11}^T & \cdots & \mathbf{A}_{s1}^T \\
    \vdots & & \vdots \\
    \mathbf{A}_{1r}^T & \cdots & \mathbf{A}_{sr}^T
  \end{pmatrix}
\end{equation*}

\subsubsection{分块对角矩阵}

设$\mathbf{A}$为$n$阶方阵，若$\mathbf{A}$的分块矩阵只有在对角线上有非零子块，其余子块都是零矩阵，且在对角线上的子块都是方阵，即
\begin{equation*}
  \mathbf{A}=\begin{pmatrix}
    \mathbf{A}_1 & & & \mathbf{O} \\
    & \mathbf{A}_2 & & \\
    & & \ddots & \\
    \mathbf{O} & & & \mathbf{A}_s
  \end{pmatrix}
\end{equation*}

其中$\mathbf{A}_i$都是方阵，那么称$\mathbf{A}$为分块对角矩阵。

分块对角矩阵的行列式有如下性质：
\begin{equation*}
  |\mathbf{A}|=|\mathbf{A}_1|\ |\mathbf{A}_2|\ \cdots |\mathbf{A}_s|.
\end{equation*}

若$|\mathbf{A}_i|$均$\neq 0$，那么$|\mathbf{A}|$也$\neq 0$，且有
\begin{equation*}
  \mathbf{A}^{-1}=\begin{pmatrix}
    \mathbf{A}_1^{-1} & & & \mathbf{O} \\
    & \mathbf{A}_2^{-1} & & \\
    & & \ddots & \\
    \mathbf{O} & & & \mathbf{A}_s^{-1}
  \end{pmatrix}
\end{equation*}

行向量默认要带一个转置符号，因为向量默认是竖写的，而行向量是横的。列向量则不用带转置符号。

\section{矩阵的初等变换与线性方程组}

\subsection{矩阵的初等变换}

\subsubsection{定义}

矩阵的初等行变换：
\begin{itemize}
  \item 对换两行
  \item 以数$k\neq 0$乘某一行中的所有元
  \item 把某一行所有元的$k$倍加到另一行对应的元上去。
\end{itemize}

初等列变换可以简单类比。

若$\mathbf{A}$经过有限次初等行变换得到$\mathbf{B}$，就说$\mathbf{A}$与$\mathbf{B}$\textbf{行等价}，记作$\mathbf{A}\overset{r}{\sim}\mathbf{B}$。同理\textbf{列等价}记作$\mathbf{A}\overset{c}{\sim}\mathbf{B}$。

矩阵的等价关系满足反身性、对称性和传递性。

非零矩阵若满足：
\begin{itemize}
  \item 非零行在零行的上面；
  \item 非零行的首非零元所在列在上一行（若有）的首非零元素所在列的右面
\end{itemize}

则称此矩阵为\textbf{行阶梯形矩阵}。若行阶梯形矩阵还满足：
\begin{itemize}
  \item 非零行的首非零元为1
  \item 首非零元所在列的其他元均为0
\end{itemize}

则称此矩阵为\textbf{行最简形矩阵}。任意矩阵总可以经过有限次初等\textbf{行}变换，变成行最简形矩阵。

对行最简形矩阵再进行初等\textbf{列}变换，使之满足条件：矩阵的左上角是一个单位矩阵，其余元素全为0。则此矩阵称为\textbf{标准形}。

由单位矩阵$\mathbf{E}$经过一次初等变换得到的矩阵称为\textbf{初等矩阵}。

\subsubsection{性质}

设$\mathbf{A}$和$\mathbf{B}$为$m\times n$矩阵，则：
\begin{itemize}
  \item $\mathbf{A}\overset{r}{\sim}\mathbf{B}$的充要条件是存在$m$阶可逆矩阵$\mathbf{P}$，使得$\mathbf{PA}=\mathbf{B}$。
  \item $\mathbf{A}\overset{c}{\sim}\mathbf{B}$的充要条件是存在$n$阶可逆矩阵$\mathbf{Q}$，使得$\mathbf{AQ}=\mathbf{B}$。
  \item $\mathbf{A}\sim\mathbf{B}$的充要条件是存在$m$阶可逆矩阵$\mathbf{P}$及$n$阶可逆矩阵$\mathbf{Q}$，使得$\mathbf{PAQ}=\mathbf{B}$。
\end{itemize}

\section{向量组的线性相关性}

\section{相似矩阵及二次型}

\section{线性空间与线性变换}

\end{document}